% Oversells «Lugʻat» shabloni (XeLaTeX yoki LuaLaTeX).
% \harf{A} — harf sarlavhasi; \soz{bosh soʻz}{turkum}{tarjima} — lugʻat maqolasi;
% \misol{...} — misol. Bet tepasida: chapda betdagi birinchi, oʻngda oxirgi
% bosh soʻz (sahifa.sty ustasi, «Sahifa» tugmasi).
\documentclass[10pt,twoside]{article}
\usepackage[a5paper]{geometry}
\usepackage{fontspec}
\setmainfont{Noto Serif}
\usepackage{babel}
\babelprovide[import=uz-Latn, main]{uzbek}
\babelprovide[import]{russian}
\def\uzbekhyphenmins{22}
\usepackage{sahifa} % Oversells sahifa dizayni (sahifa.sty)

% Every headword is a mark: the head shows the first (\rightmark) and the last (\leftmark) of the page.
\NewDocumentCommand\soz{mmmo}{\begin{uslub}{maqola}\markboth{#1}{#1}\matnuslubi{boshsoz}{#1} \matnuslubi{turkum}{#2}
  \foreignlanguage{russian}{#3}\IfValueT{#4}{; #4}.\end{uslub}}
\newcommand\misol[1]{\matnuslubi{misol}{#1}}
\newcommand\harf[1]{\section*{#1}}

\begin{document}

\begin{ustunlar}{2}
\harf{A}
\soz{ajoyib}{sft}{удивительный, чудесный}[\misol{ajoyib kun} — чудесный день]
\soz{aka}{ot}{старший брат}
\soz{aql}{ot}{ум, разум}[\misol{aql bilan ish tutmoq} — поступать разумно]
\soz{arava}{ot}{телега, повозка}
\soz{asal}{ot}{мёд}[\misol{asal ari} — пчела]
\soz{ayol}{ot}{женщина}
\soz{aytmoq}{feʼl}{говорить, сказать}[\misol{rost aytmoq} — говорить правду]

\harf{B}
\soz{bahor}{ot}{весна}[\misol{erta bahor} — ранняя весна]
\soz{baland}{sft}{высокий}
\soz{bayram}{ot}{праздник}[\misol{bayram tabrigi} — праздничное поздравление]
\soz{bilim}{ot}{знание}
\soz{bogʻ}{ot}{сад}[\misol{olma bogʻi} — яблоневый сад]
\soz{bola}{ot}{ребёнок}
\soz{bormoq}{feʼl}{идти, ехать}[\misol{uyga bormoq} — идти домой]
\soz{bulut}{ot}{облако, туча}

\harf{D}
\soz{daftar}{ot}{тетрадь}
\soz{daraxt}{ot}{дерево}[\misol{mevali daraxt} — плодовое дерево]
\soz{daryo}{ot}{река}
\soz{dasturxon}{ot}{скатерть; угощение}[\misol{dasturxon yozmoq} — накрывать на стол]
\soz{doʻst}{ot}{друг}[\misol{eski doʻst} — старый друг]
\soz{dunyo}{ot}{мир, свет}

\harf{G}
\soz{gul}{ot}{цветок}[\misol{gul bogʻi} — цветник]
\soz{goʻzal}{sft}{красивый}
\soz{gapirmoq}{feʼl}{говорить, разговаривать}

\harf{O}
\soz{ona}{ot}{мать}[\misol{ona tili} — родной язык]
\soz{osmon}{ot}{небо}
\soz{oʻqimoq}{feʼl}{читать; учиться}[\misol{kitob oʻqimoq} — читать книгу]
\soz{oʻrik}{ot}{абрикос}
\soz{ovqat}{ot}{еда, пища}

\end{ustunlar}

\end{document}
